\section*{Five principles}\label{sec:organizing-principles}
\addcontentsline{toc}{section}{Five principles}

\subsection*{Principle 1 (Matrix-Probe Measurement)}

\emph{A stack of \(N\) Dirac branes on the topological line is a
finite Chan--Paton representation
\(\rho_\phi\colon A=\C[[z_1,z_2]]\to\End(\C^N)\),
\(z_i\mapsto\phi_i\), forced onto the commuting variety
\([\phi_1,\phi_2]=0\) by the conjugation moment map; a Hamiltonian
\(f\) on the formal symplectic normal polydisk is measured on the
brane as the gauge-invariant trace character of its representation,
\(J(f)=\overline{\operatorname{Tr}\,\rho_\phi(f)}=
\overline{\operatorname{Tr}\,f(\phi_1,\phi_2)}\).}

The substitution \(z_i\rightsquigarrow\phi_i\) into a pair of
\(\mathfrak{gl}_N\)-valued matrices is the only datum.  The
\(GL_N\)-symplectic potential \(\operatorname{Tr}(\phi_1\,d\phi_2)\),
the moment map \(-\operatorname{Tr}(\epsilon[\phi_1,\phi_2])\), the
Dirac constraint \([\phi_1,\phi_2]=0\), and the Koszul derived
zero-fibre model \(Q\psi=[\phi_1,\phi_2]\) follow.  The brane
represents the commutative coordinate ring \(A\); the Hamiltonian
Lie algebra \(\mathfrak h=A/\C\!\cdot\!1\) does not act on the Chan--Paton
space \(\C^N\) by commuting operators, but on the representation
moduli \(\Rep_N(A)/\!/GL_N\) by Hamiltonian vector fields, and the
closed-to-open measurement is the trace character above.  Cyclic
invariance and adjacent-swap \(Q\)-exactness collapse every trace
word to the symmetric symbol family
\(\operatorname{Tr}(\phi_1^k\phi_2^l)=J(z_1^kz_2^l)\).  After formal
Darboux normalization at any brane point of any connected
holomorphic-symplectic surface, every admissible protected
Hamiltonian sector reduces to the same brane microscope, and
divergence is read coordinate by coordinate in the four-class
residual \(\operatorname{Ob}^{\mathrm{univ}}_{p,X}\).

\subsection*{Principle 2 (The Centre of Mass Remembers What Hamiltonian Reduction Forgets)}

\emph{The constant Hamiltonian generates the zero Hamiltonian vector
field closed-side and survives open-side as the brane \(U(1)\)
centre of mass; the same scalar cocycle is detected algebraically as
\([\bar c]\in H^2_{\mathrm{Lie}}(\bar A,\C)\), Lie-theoretically as the
central extension of \(\bar A\), in the trace algebra as
\(\{\operatorname{Tr}\phi_1,\operatorname{Tr}\phi_2\}=N\), and after
Moyal quantization as the Capelli/QME shift \(\hbar N[\bar c]\).}

The four detections are not four cocycles -- they are one cocycle
transgressed across four faces of the same scalar-reduction step.
Removing constants from \(\C[[z_1,z_2]]\) to obtain \(\mathfrak h\)
projects out the kernel of the Hamiltonian vector field map
\(f\mapsto X_f\); on the brane the same kernel reappears as the
centre-of-mass mode, so the closed-open mismatch is forced and equal
to the transgressed class.  Algebraic deformation quantization of the
polydisk by the Moyal star product converts the same scalar into the
order-\(\hbar\) Capelli term, with \(\hbar N\) the only admissible
coefficient; the first scalar-reduced \(\mathfrak{gl}_N\) QME
obstruction is therefore \(W_1^{\mathrm{sc,ord}}=\hbar N\,[\bar c]\),
and any cancellation must run through the unreduced central-character
branch, the opposite-extension branch, or the balanced-supertrace
branch.

\subsection*{Principle 3 (Shifted Cotangent and Residue Pairing)}

\emph{The closed BV sector is the shifted-cotangent local Lie algebra
\(\mathfrak g=\mathfrak h\ltimes\mathfrak h^\vee_{\mathrm{cont}}[1]\).
A Hamiltonian \(f\in\mathfrak h\) has two CE coordinate avatars on
\(\mathfrak g\): the Hamiltonian-dual \(c_f\) acts on observables as
the Schouten field \(\theta_f\); the cotangent-dual \(u_f\) is the
closed insertion the brane reads, with image \(O_f=J(f)\).  The
continuous dual \(\mathfrak h^\vee_{\mathrm{cont}}\) is realized by
the residue pairing as the Grothendieck--Matlis principal-part module
\(\mathcal P=H^2_{\mathfrak m}(A)\otimes\omega_A
=\bigoplus_{a+b>0}\C\,\rho_{a,b}\) --- the holomorphic-surface analog
of the Heisenberg current.}

The CE/PV cochain dictionary on generators is
\[
  c^I\longmapsto\theta^I,\qquad u_I\longmapsto O_I .
\]
The Hamiltonian-dual coordinate \(c_f\) is a perfect Lie cochain
acting on observables as the Schouten field \(\theta_f\); a perfect
Hamiltonian Lie algebra carries no closed length-one CE cochain in
the \(c\)-direction, so \(c_f\) acts on the observable algebra
without itself entering it as an observable.  The
shifted-cotangent-dual coordinate \(u_f\) is the closed insertion the
brane reads, with \(\theta_f(O_g)=O_{\{f,g\}}\); the brane functional
\(J\) implements the duality, and the reduced CE/PV
central-operation theorem links the two avatars.

The brane defect line carries a canonical polarization: configuration
observables are polynomial Taylor classes
\(O_f=\overline{\operatorname{Tr}f(\phi_1,\phi_2)}\) on the
matrix-valued normal coordinate; cotangent defect momenta are
distributional residues
\(\rho_{a,b}=z_1^{-a-1}z_2^{-b-1}\,dz_1dz_2\) in
\(\mathcal P=\bigoplus_{a+b>0}\C\,\rho_{a,b}\).  The residue pairing
makes the identification \(\mathfrak h^\vee_{\mathrm{cont}}\cong
\mathcal P\) canonical, and the residue-coadjoint formula
\(z_1^p z_2^q\cdot\rho_{a,b}=-(pb-qa+p-q)\,\rho_{a-p+1,b-q+1}\) is
its closed-form fingerprint.  The polynomial PBW one-\(\psi\)
descendants \(\Psi_{a,b}\) carry the same labels and a distinct
\(\mathfrak h\)-module structure: linear Hamiltonians act locally
nilpotently on the polynomial side and pole-raisingly on
\(\mathcal P\); the PBW classes are the special-fibre shadow, the
deformation-level cotangent module is \(\mathcal P\), and the
Fourier-polarized Rees bridge is the unique passage between them.
Boundary Hamiltonian observables are the image of the cotangent
coordinate \(u\).

\subsection*{Principle 4 (Three Independent Localities)}

\emph{The factorization-algebra structure on the brane line is forced
by three independent local-current identities, with \((t,t')\) two
brane-time points and \(s\) reserved for the transverse \(\R_s\)
coordinate of \(X=\R_t\times\R_s\times\C^2\),
\[
  \delta(t-t')\quad\text{(support)},\qquad
  L_v=[Q,\iota_v]\quad\text{(topological)},\qquad
  L_{\partial/\partial\bar z_i}
   =[Q,\iota_{\partial/\partial\bar z_i}]
   \quad\text{(formal-holomorphic)},
\]
and none of the three implies the others.}

Support locality identifies \(P_0\)-brackets of disjoint smeared
Hamiltonian currents with zero; topological locality makes the brane
factorization algebra locally constant in \(\R^2_{\mathrm{top}}\);
formal-holomorphic locality concentrates closed-side cohomology on
holomorphic jets at the brane point and pairs cotangent modes with
Taylor Hamiltonians by residue.  Trace observables exhibit all three
identities; shifted-cotangent currents exhibit support and topological
locality but \emph{not} formal-holomorphic locality, and the gap is
precisely the cotangent centrality homotopy that no single identity
supplies.  The unique factorization algebra compatible with the three
identities is the closed mixed-HT current factorization algebra
\(\mathcal F^{\mathrm{cl}}_{\mathrm{Ham}}(O)
=C^{\mathrm{CE}}_\bullet\bigl(\Omega_c^\bullet(U)\widehat\otimes
\Omega_c^{0,\bullet}(B)\widehat\otimes\mathfrak g\bigr)\) on
\(O=U\times B\subset\R^2_{\mathrm{top}}\times\C^2_{\mathrm{hol}}\),
\(\mathfrak g=\mathfrak h\ltimes\mathfrak h^\vee_{\mathrm{cont}}[1]\)
(Theorem~\ref{thm:mixed-HT-factorization-theorem}).  Local
constancy on a topological disk gives the bulk \(E_2\)-algebra
\(\mathcal A^{\mathrm{cl}}_{\mathrm{bulk}}(B)
=\mathcal F^{\mathrm{cl}}_{\mathrm{Ham}}(D^2\times B)\)
(Corollary~\ref{cor:E2-bulk-algebra}); local constancy on the
brane worldline gives the brane \(E_1\)-algebra
\(A^{\mathrm{cl}}_{\partial,N}=
C^\bullet_{\mathrm{CE}}(\mathfrak{gl}_N,\mathrm{Kosz}_N)\)
(Definition~\ref{def:E1-brane-algebra}); the closed-to-open
derived-centre map
\(\Phi_N\colon\mathcal A^{\mathrm{cl}}_{\mathrm{bulk},0}\to
Z^{\mathrm{der}}_{E_1}(A^{\mathrm{cl}}_{\partial,N})\) has
reduced-coordinate shadow the CE/PV dictionary
(Theorem~\ref{thm:closed-to-open-derived-centre-map}).  The native
bulk object is a holomorphic factorization algebra on \(\C^2\)
valued in \(E_2\)-algebras, not a curve vertex algebra.

\subsection*{Principle 5 (Obstruction Classes are Curvature Coordinates for Stronger Targets)}

\emph{Every stronger global, quantum, chain-level, or all-bidegree
construction is a filtered \(L_\infty\) transport of the local
matrix-microscope Maurer--Cartan datum \(\Theta_{\mathrm{loc}}\) into
a target deformation complex \(\mathfrak K_{\mathcal T}\), with master
equation \(F_{\mathcal T}(\Theta_{\mathrm{loc}}+\tau)=0\).  The four
named curvature classes
\(\Omega^{\mathrm{rad}}_{a,b}\),
\(\operatorname{Ob}_{\mathrm{cent/QME/desc}}\),
\([\operatorname{Ob}^{\mathrm{red}}_{w,\partial,\hbar}]\),
\(\operatorname{Ob}^{\Pi}_{\mathrm{BM}}\)
are the curvature components of the local-to-global transport:
zero curvature gives strict coherence; \(d_\Theta\)-exact curvature
gives counterterm or homotopy coherence; central, closed,
non-exact curvature gives projective or gerbe-valued coherence;
non-absorptive curvature forces an enlargement of the target along a
substitution \(\mathcal T\rightsquigarrow\mathcal T'\).}

The instance \(\mathcal T=\mathrm{Costello\ QME}\) is the
quantization compatibility test.  The bidegree square
\[
  f\xrightarrow{\;J\;}J(f)\qquad
  f\xrightarrow{\;\mathrm{Weyl}\;}\Phi_\hbar(f)
   \xrightarrow{\;\overline{\operatorname{Tr}}\;}
   J_\hbar(f)\coloneq\overline{\operatorname{Tr}\,\Phi_\hbar(f)}
\]
asks whether Weyl ordering and the reduced trace
\(\overline{\operatorname{Tr}}\) commute up to controlled
counterterm.  The verified range
(Theorem~\ref{thm:componentwise-quantum-coefficient-comparison})
covers the balanced diagonal through \((6,6)\), all non-edge
bidegrees through total degree \(13\), selected total-degree-\(14\)
cases, and one-moving-letter edge strips; outside it the obstruction
is the per-bidegree class \(\Omega^{\mathrm{rad}}_{a,b}\),
equivalently the decorated Stokes failure for
\(D^\square_{a,b}=C^+_{a,b}\partial_2\).  At order \(\hbar^1\), the
scalar shift \(\hbar N[\bar c]\) of Principle~2 is the universal
Capelli term --- the trace-\(U(1)\) cocycle of \(N\) coincident
D-branes on noncommutative \(\C^2_\hbar\); at higher graph order, the
controlling class is
\([\operatorname{Ob}^{\mathrm{red}}_{w,\partial,\hbar}]\in
H^1\bigl(\mathfrak Q^\bullet_{w,\partial,\hbar}(I),\,d_{\mathrm{QME}}\bigr)\)
in the filtered weighted brane-defect local-functional complex,
with the order-three CE-source row \(\theta_3\) finite-window
prototype:
\(H^1_{\mathrm{bare}}=\C\!\cdot\!\mathsf E_{\theta_3,M}\), killed by
the unique counterterm $d_M C_{\theta_3,M}=-\mathsf E_{\theta_3,M}$,
\(c=1\), in any finite window $M\ni\theta_3$.

A target \(\mathcal T\) -- global Hamiltonian descent, unreduced
\(P_0\)-factorization-centre lift, Costello QME quantization,
all-bidegree Weyl/Moyal radial normal form, strict all-window
Bochner--Martinelli current transfer, modular open-closed clutching,
CY-to-chiral specialisation -- carries a filtered dg Lie deformation
complex \(\mathfrak K_{\mathcal T}\) controlling its coherence, and
the local theorem produces a transported Maurer--Cartan datum
\[
  \Theta_{\mathrm{loc}}\in F^1\mathfrak K^1_{\mathcal T} .
\]
The strict \(\mathcal T\)-realisation asks
\(F_{\mathcal T}(\Theta_{\mathrm{loc}})=0\); a twisted realisation is a
filtered cochain \(\tau\) (or central extension of
\(\mathfrak K_{\mathcal T}\)) with
\[
  F_{\mathcal T}(\Theta_{\mathrm{loc}}+\tau)=0
  \qquad\text{in the extended target.}
\]
The obstruction class \([F_{\mathcal T}(\Theta_{\mathrm{loc}})]\)
classifies the first possible twist of the upgrade: the scalar anomaly
\(\hbar N[\bar c]\) is the central-extension twist that promotes the
strict scalar-reduced action to the projective Hamiltonian action; the
reduced non-scalar QME class
\([\operatorname{Ob}^{\mathrm{red}}_{w,\partial,\hbar}]\) is the
counterterm twist of the Costello quantum master equation, with the
\(\theta_3\) row exhibiting the prototype
\((dC=-E,\,c=1)\) cancellation; the four-component vector
\(\operatorname{Ob}_{\mathrm{cent/QME/desc}}\) is the homotopy-centre
twist passing the reduced cochain CE/PV map to the derived
\(P_0\)-factorization centre \(Z^{P_0,\mathrm{der}}_{\mathrm{fact}}\);
the radial obstruction
\(\Omega^{\mathrm{rad}}_{a,b}\in\operatorname{coker}B_{a,b}\) is the
decorated radial Stokes twist enlarging the strict Weyl normal form
to the decorated PBW/Stokes target; the strict Bochner--Martinelli
class \(\operatorname{Ob}^{\Pi}_{\mathrm{BM}}\) is nonzero in every
fixed finite current window; the Hamiltonian factorization current
acts coherently on the pro-Matlis tower
\(\widehat{\mathcal P}_\Pi=\varprojlim_N\mathcal P_{\le N}\);
the Hamiltonian period \(\mathrm{per}_X(v)=[h_i-h_j]\in H^1(X,\C_X)\)
is the descent twist promoting strict global Hamiltonians to sections
of a \(\C^\times\)-Hamiltonian gerbe.  Some classes are integrable
twists inside the original target; others are no-go selectors that
mark a substitution \(\mathcal T\rightsquigarrow\mathcal T'\) as
the next step.  The obstruction calculus is simultaneously the
boundary of what the local theorem proves and the deformation complex
of possible global quantum completions.  The master schema is
Theorem~\ref{thm:obstruction-twist-dichotomy} of
Appendix~\ref{app:master-deformation-complex}; every named obstruction
class is a clause of that schema, and stronger constructions live in
the spaces \(\mathrm{TwGlob}_{\mathcal T}(\Theta_{\mathrm{loc}})\) of
admissible twists.

Every frontier obstruction is a named cohomology class or a named
contracting-homotopy obligation; the corresponding repair is either
the smallest enlargement of the ambient algebra in which the intended
mechanism is a theorem, or the proved obstruction theorem naming the
exact class that prevents it.
